\clearpage
\gdef\otChapNum{0}\pagestyle{otFront}
\phantomsection\label{defis-debut}\pdfbookmark[0]{Défis d'approfondissement}{bm-defis}
\otPageTitre{Défis d'approfondissement}{Chercher plus loin avec les outils de Première}
\fontsize{11}{16}\selectfont
Les \otNbExosChapitres\ exercices des chapitres installent les connaissances et les méthodes.
Les \otNbDefis\ défis suivants, un par chapitre, portent le total à \textbf{\otNbExosTotal\ exercices corrigés}.
Ils prolongent chaque thème par une question qui demande de prendre du recul, d'estimer un ordre de grandeur,
de construire un modèle ou de relier plusieurs notions.\par\vspace{4mm}
\textbf{Avancé} désigne une modélisation ou une synthèse soutenue.
\textbf{Très exigeant} ajoute une discussion, une démonstration ou plusieurs changements de point de vue.
Ces indications sont des repères pédagogiques, pas un classement officiel.\par\vspace{4mm}
Les défis sont facultatifs pour un premier passage. Ils restent dans le cadre du programme ; un outil qui
le dépasserait est fourni dans l'énoncé. Les durées proposées sont indicatives.\par\vspace{5mm}
\editionEtape{Chercher avant de lire la solution}
Faire un schéma, écrire les données et les relations utiles, tester un cas simple, puis noter ce qui manque
pour conclure. En cas de blocage, lire seulement le premier indice, puis reprendre la recherche.
Les indications sont séparées des énoncés et des corrigés.
\begin{center}\renewcommand{\arraystretch}{1.25}
\begin{tabular}{@{}lll@{}}\toprule
Défis & Thèmes & Préparation conseillée\\\midrule
D01 à D10 & Constitution et transformations de la matière & Chapitres 1 à 10\\
D11 à D13 & Mouvement et interactions & Chapitres 11 à 13\\
D14 à D16 & L'énergie : conversions et transferts & Chapitres 14 à 16\\
D17 à D20 & Ondes et signaux & Chapitres 17 à 20\\
D21 & Mesures et incertitudes & Chapitre 21\\\bottomrule
\end{tabular}\end{center}
\editionRetenir{Une solution incomplète mais justifiée est un progrès.
Après la correction, refaire la question clé sans regarder, puis expliquer pourquoi la méthode fonctionne.}

\editionPage{defi-ex-1}{D01 · Le dernier souffle de César}{Recherche et rédaction}{Avancé · 30 à 45 min · Chapitre 1\hfill\hyperref[defi-cor-1-1]{Corrigé p.~\pageref*{defi-cor-1-1}}}
\begin{blocEdition}{defi-ex-1}
En 44 avant notre ère, Jules César rend son dernier souffle. Un problème classique d'estimation pose la question
suivante : combien de molécules de ce dernier souffle respirons-nous aujourd'hui, à chaque inspiration ?\par\smallskip
On admet qu'en plus de deux mille ans, les molécules de ce souffle se sont réparties uniformément dans toute
l'atmosphère, et qu'elles n'ont été ni détruites ni retirées de l'air.
\begin{donnees}
Masse de l'atmosphère terrestre : \qty{5,1e18}{kg}. Masse molaire moyenne de l'air : \qty{29}{g.mol^{-1}}. Volume d'air d'une
inspiration (ou d'une expiration) au repos : \qty{0,50}{L}, à \qty{20}{\degreeCelsius} et \qty{1013}{hPa}, où $V_\mathrm{m}=\qty{24,0}{L.mol^{-1}}$.
$N_\mathrm{A}=\qty{6,02e23}{mol^{-1}}$.
\end{donnees}
\begin{questions}
\item Calculer la quantité de matière d'air de l'atmosphère, puis le nombre de molécules qu'elle contient.
\item Calculer la quantité de matière d'air d'une inspiration, puis le nombre de molécules qu'elle contient.
\item Combien d'inspirations l'atmosphère contient-elle ? Comparer ce nombre au nombre de molécules d'une inspiration.
\item Quelle fraction de l'atmosphère le dernier souffle de César représente-t-il ? En déduire le nombre moyen de
molécules de ce dernier souffle contenues dans chacune de nos inspirations.
\item Au repos, on respire environ quinze fois par minute. Combien de molécules du dernier souffle de César
respire-t-on en une heure ?
\item Discuter les hypothèses : le dioxygène est consommé par les êtres vivants et renouvelé par la photosynthèse, le
diazote représente 78~\% de l'air et réagit très peu, une partie du dioxyde de carbone se dissout dans les océans.
L'ordre de grandeur du résultat est-il remis en cause ?
\end{questions}
\end{blocEdition}

\editionPage{defi-ex-2}{D02 · Deux colorants, une seule couleur}{Recherche et rédaction}{Avancé · 40 à 50 min · Chapitre 2\hfill\hyperref[defi-cor-2-1]{Corrigé p.~\pageref*{defi-cor-2-1}}}
\begin{blocEdition}{defi-ex-2}
La couleur verte d'un sirop de menthe vient de deux colorants : la tartrazine, notée J, jaune, et le bleu patenté,
noté B. On a mesuré l'absorbance du sirop dilué vingt fois à deux longueurs d'onde, dans une cuve de \qty{1,00}{cm} :
$A_{427}=0{,}62$ à \qty{427}{nm} et $A_{638}=0{,}53$ à \qty{638}{nm}. Lorsque plusieurs espèces absorbent, leurs absorbances s'ajoutent :
$A=k_\mathrm{J}\,c_\mathrm{J}+k_\mathrm{B}\,c_\mathrm{B}$, chaque coefficient dépendant de la longueur d'onde.
\begin{center}
\begin{tabular}{lcc}
\toprule
Coefficient $k$ (\unit{L.\micro mol^{-1}}) & à \qty{427}{nm} & à \qty{638}{nm} \\
\midrule
tartrazine (J) & 0,0270 & 0 \\
bleu patenté (B) & 0,0060 & 0,106 \\
\bottomrule
\end{tabular}
\end{center}
\begin{questions}
\item Pourquoi ne peut-on pas déterminer la concentration de la tartrazine à partir de la seule mesure à \qty{427}{nm} ?
\item Écrire les expressions de $A_{427}$ et de $A_{638}$ en fonction des concentrations $c_\mathrm{J}$ et $c_\mathrm{B}$ de la solution diluée.
\item En déduire $c_\mathrm{B}$, puis $c_\mathrm{J}$.
\item Quelle erreur relative commettrait-on sur $c_\mathrm{J}$ en oubliant l'absorption du bleu patenté à \qty{427}{nm} ?
\item Calculer les concentrations des deux colorants dans le sirop.
\item Dans le cas général, les deux colorants absorbent aux deux longueurs d'onde. Écrire le système d'équations obtenu et
expliquer comment le résoudre. Pourquoi choisit-on deux longueurs d'onde où les deux colorants absorbent de façons très
différentes ?
\end{questions}
\end{blocEdition}

\editionPage{defi-ex-3}{D03 · L'arbre de Diane prend du poids}{Recherche et rédaction}{Avancé · 40 à 50 min · Chapitre 3\hfill\hyperref[defi-cor-3-1]{Corrigé p.~\pageref*{defi-cor-3-1}}}
\begin{blocEdition}{defi-ex-3}
Un fil de cuivre de \qty{5,00}{g}, enroulé en spirale, est plongé dans \qty{100}{mL} d'une solution de nitrate d'argent,
$\ce{Ag+}+\ce{NO3-}$, de concentration \qty{0,100}{mol.L^{-1}}. Des aiguilles d'argent se forment sur le fil et la solution bleuit. De
temps en temps, on sort le fil, on le rince, on le sèche et on le pèse : sa masse augmente. On admet que tout l'argent
formé reste accroché au fil.
\begin{center}
\begin{tabular}{lcc}
\toprule
Élément & cuivre & argent \\
\midrule
Masse molaire (\unit{g.mol^{-1}}) & 63,5 & 107,9 \\
\bottomrule
\end{tabular}
\end{center}
\begin{questions}
\item Établir l'équation de la réaction, les couples étant $\ce{Ag+}/\ce{Ag}$ et $\ce{Cu^{2+}}/\ce{Cu}$.
\item Calculer les quantités initiales d'ions argent et de cuivre. Quel réactif sera entièrement consommé ?
\item Calculer la masse de cuivre consommée et la masse d'argent déposée lorsque la réaction s'arrête. En déduire la masse
finale du fil.
\item On note $x$ la quantité de cuivre ayant réagi à un instant donné. Montrer que la masse du fil s'écrit
$m=5{,}00+152{,}3\,x$, avec $m$ en grammes et $x$ en moles. Que représente le nombre 152,3 ?
\item Au cours de l'expérience, le fil pèse \qty{5,38}{g}. En déduire $x$, puis la concentration des ions cuivre(II) dans la
solution à cet instant.
\item Quel volume de cette solution de nitrate d'argent faudrait-il pour dissoudre entièrement le fil ?
\item Proposer une méthode pour suivre l'avancement de la réaction sans sortir le fil de la solution.
\end{questions}
\end{blocEdition}

\editionPage{defi-ex-4}{D04 · La sonde lambda}{Recherche et rédaction}{Avancé · 45 à 60 min · Chapitre 4\hfill\hyperref[defi-cor-4-1]{Corrigé p.~\pageref*{defi-cor-4-1}}}
\begin{blocEdition}{defi-ex-4}
Dans un moteur à essence, le carburant brûle dans l'air aspiré par le moteur. Une sonde placée dans le pot
d'échappement, la sonde lambda, mesure le dioxygène restant ; le calculateur ajuste alors la quantité d'essence injectée
pour maintenir le rapport
\[\lambda=\frac{\text{quantité de dioxygène introduite}}{\text{quantité de dioxygène d'un mélange stœchiométrique}}\]
très proche de 1. On assimile l'essence à l'octane $\ce{C8H18}$, et l'air à un mélange contenant 21~\% de dioxygène en quantité de
matière, de masse molaire moyenne \qty{29,0}{g.mol^{-1}}.
\begin{center}
\begin{tabular}{lccc}
\toprule
Élément & H & C & O \\
\midrule
Masse molaire (\unit{g.mol^{-1}}) & 1,0 & 12,0 & 16,0 \\
\bottomrule
\end{tabular}
\end{center}
\begin{questions}
\item Établir l'équation de la combustion complète de l'octane, qui produit du dioxyde de carbone et de l'eau.
\item Pour \qty{1,0}{mol} d'octane, quelle quantité de dioxygène forme un mélange stœchiométrique ? En déduire la quantité, puis
la masse d'air correspondante.
\item Montrer que le rapport massique air/essence du mélange stœchiométrique vaut environ 15,1. Les motoristes retiennent
14,7 pour l'essence réelle : proposer une explication.
\item Pour \qty{1,0}{mol} d'octane, déterminer le réactif limitant et la composition de l'état final pour $\lambda=0{,}90$, puis pour
$\lambda=1{,}10$. On pourra utiliser le programme de l'essentiel.
\item Un mélange « riche » ($\lambda<1$) rejette des imbrûlés et du monoxyde de carbone. Expliquer ce constat à l'aide de la
question précédente.
\item Une voiture consomme \qty{6,0}{L} d'essence aux \qty{100}{km} ; la masse volumique de l'essence vaut \qty{0,75}{kg.L^{-1}}. Calculer la
masse de dioxyde de carbone émise par kilomètre.
\end{questions}
\end{blocEdition}

\editionPage{defi-ex-5}{D05 · Que vaut vraiment une eau de Javel ?}{Recherche et rédaction}{Avancé · 45 à 60 min · Chapitre 5\hfill\hyperref[defi-cor-5-1]{Corrigé p.~\pageref*{defi-cor-5-1}}}
\begin{blocEdition}{defi-ex-5}
Une eau de Javel « 2,6~\% de chlore actif » libère, en milieu acide, \qty{2,6}{g} de dichlore pour \qty{100}{mL}. Ses ions
hypochlorite $\ce{ClO-}$ ne réagissent pas assez vite avec un réactif titrant usuel : on procède à un \textbf{titrage indirect}.
\par\smallskip\noindent\textbf{Protocole.} On dilue l'eau de Javel dix fois. On prélève \qty{10,0}{mL} de solution diluée, on ajoute environ \qty{1}{g}
d'iodure de potassium $\ce{KI}$ (excès) et de l'acide éthanoïque : les ions hypochlorite oxydent les ions iodure en diiode,
$\ce{ClO- + 2I- + 2H+ -> I2 + Cl- + H2O}$. On titre ensuite le diiode formé par du thiosulfate de sodium à \qty{5,00e-2}{mol.L^{-1}}.
Pour un flacon neuf, $V_\mathrm{E}=\qty{14,6}{mL}$ ; pour un flacon ouvert depuis un an et rangé au soleil, $V_\mathrm{E}=\qty{11,2}{mL}$.
\begin{center}
\begin{tabular}{lccc}
\toprule
Donnée & $M(\ce{Cl})$ & $M(\ce{KI})$ & Équation du titrage \\
\midrule
Valeur & \qty{35,5}{g.mol^{-1}} & \qty{166,0}{g.mol^{-1}} & $\ce{I2 + 2S2O3^{2-} -> 2I- + S4O6^{2-}}$ \\
\bottomrule
\end{tabular}
\end{center}
\begin{questions}
\item Vérifier que l'équation de la réaction entre les ions hypochlorite et iodure est ajustée. Quel est l'oxydant ? le réducteur ?
\item Pourquoi faut-il mettre les ions iodure en excès ? Vérifier que \qty{1}{g} d'iodure de potassium suffit.
\item Relier la quantité de thiosulfate versée à l'équivalence à la quantité de diiode, puis à celle des ions hypochlorite.
\item Calculer la concentration en ions hypochlorite de l'eau de Javel du flacon neuf.
\item En milieu acide, $\ce{ClO- + Cl- + 2H+ -> Cl2 + H2O}$. Calculer le « pourcentage de chlore actif » du flacon neuf et le comparer à
l'étiquette.
\item Quelle part de son pouvoir désinfectant le vieux flacon a-t-il perdue ? Quel conseil de conservation en tirer ?
\end{questions}
\end{blocEdition}

\editionPage{defi-ex-6}{D06 · La molécule d'eau en géométrie}{Recherche et rédaction}{Très exigeant · 50 à 60 min · Chapitre 6\hfill\hyperref[defi-cor-6-1]{Corrigé p.~\pageref*{defi-cor-6-1}}}
\begin{blocEdition}{defi-ex-6}
\textbf{Partie A : l'angle du tétraèdre.} On place l'atome de carbone du méthane au centre d'un cube d'arête 2 (unité
arbitraire), et les quatre atomes d'hydrogène sur quatre sommets du cube qui ne sont pas reliés par une arête : deux
atomes d'hydrogène sont toujours aux extrémités d'une diagonale d'une face.
\begin{questions}
\item Montrer que la distance C--H vaut $\sqrt3$ et que la distance entre deux atomes d'hydrogène vaut $2\sqrt2$.
\item À l'aide de la formule d'Al-Kashi, $a^2=b^2+c^2-2bc\cos\theta$, montrer que l'angle $\theta$ entre deux liaisons C--H vérifie
$\cos\theta=-\dfrac13$, puis calculer $\theta$.
\item Dans l'ammoniac et dans l'eau, les angles valent \qty{107}{\degree} et \qty{104,5}{\degree}. Proposer une explication.
\end{questions}
\textbf{Partie B : mesurer la polarité.} Les chimistes mesurent la polarité d'une molécule par son moment dipolaire, en
debyes (D). On représente chaque liaison O--H par un vecteur $\vec p_1$ de norme $p_1=\qty{1,51}{D}$, dirigé de l'atome d'hydrogène
vers l'atome d'oxygène ; le moment dipolaire de la molécule est la norme de la somme des vecteurs de ses liaisons.
\begin{questions}[resume]
\item Montrer que, pour un angle $\alpha$ entre les liaisons, la norme de la somme vaut $p=2\,p_1\cos\left(\dfrac{\alpha}{2}\right)$.
\item Calculer $p$ pour la molécule d'eau. La valeur mesurée est \qty{1,85}{D} : conclure.
\item Que vaudrait $p$ si la molécule d'eau était linéaire ? si son angle valait \qty{109,5}{\degree} ? Que vaut le moment dipolaire du
dioxyde de carbone ? Relier ces résultats à la notion de molécule polaire.
\end{questions}
\end{blocEdition}

\editionPage{defi-ex-7}{D07 · L'étiquette d'une eau minérale}{Recherche et rédaction}{Avancé · 45 à 60 min · Chapitre 7\hfill\hyperref[defi-cor-7-1]{Corrigé p.~\pageref*{defi-cor-7-1}}}
\begin{blocEdition}{defi-ex-7}
L'étiquette d'une eau minérale très minéralisée indique sa composition en ions, en milligrammes par litre. Une solution est
toujours électriquement neutre : la quantité totale de charges positives apportées par les cations doit être égale à celle
des charges négatives apportées par les anions.
\begin{center}
\begin{tabular}{lcccc@{\qquad}lcccc}
\toprule
Cation & $\ce{Ca^{2+}}$ & $\ce{Mg^{2+}}$ & $\ce{Na+}$ & $\ce{K+}$ & Anion & $\ce{SO4^{2-}}$ & $\ce{HCO3-}$ & $\ce{Cl-}$ & $\ce{NO3-}$ \\
\midrule
\unit{mg.L^{-1}} & 486 & 84 & 14,2 & 4,1 & \unit{mg.L^{-1}} & 1187 & 403 & 14 & 4,3 \\
\bottomrule
\end{tabular}
\end{center}
Masses molaires (\unit{g.mol^{-1}}) : H 1,0 ; C 12,0 ; N 14,0 ; O 16,0 ; Na 23,0 ; Mg 24,3 ; S 32,1 ; Cl 35,5 ; K 39,1 ; Ca 40,1.
\begin{questions}
\item Calculer la masse molaire de chaque ion polyatomique, puis la concentration en quantité de matière de chaque ion,
en \unit{mmol.L^{-1}}.
\item Un ion de charge $z$ (en charges élémentaires) et de concentration $c$ apporte $|z|\,c$ « moles de charges » par litre.
Calculer la quantité de charges positives et de charges négatives apportées par litre.
\item L'étiquette est-elle cohérente avec l'électroneutralité ? Calculer l'écart relatif.
\item Comparer les concentrations des ions calcium et sulfate. Quelle roche, traversée par l'eau, peut expliquer cette
composition ? On donne : le gypse est un sulfate de calcium hydraté, $\ce{CaSO4}$, $2\,\ce{H2O}$.
\item Quelle masse de gypse un litre de cette eau a-t-il dissous, au minimum ?
\item Écrire un programme Python qui calcule et compare les deux quantités de charges à partir de dictionnaires.
\end{questions}
\end{blocEdition}

\editionPage{defi-ex-8}{D08 · L'odeur du chocolat}{Recherche et rédaction}{Avancé · 40 à 50 min · Chapitre 8\hfill\hyperref[defi-cor-8-1]{Corrigé p.~\pageref*{defi-cor-8-1}}}
\begin{blocEdition}{defi-ex-8}
Lors de la torréfaction des fèves de cacao se forment des molécules odorantes. L'une d'elles, à l'odeur de malt et de
chocolat, a pour formule brute $\ce{C5H10O}$ ; son squelette carboné est saturé et elle contient un groupe carbonyle. Trois
indices permettent de l'identifier :
\begin{questions}
\item[] \textbf{Indice A.} Son spectre infrarouge est reproduit ci-dessous.
\item[] \textbf{Indice B.} Son squelette carboné est ramifié.
\item[] \textbf{Indice C.} L'atome de carbone qui porte la ramification n'est pas lié à l'atome de carbone du groupe carbonyle.
\end{questions}
\label{figedition-defi-ex-8-1}\schema{defi-08}
\begin{questions}
\item Montrer qu'un composé carbonylé saturé de formule $\ce{C5H10O}$ est un aldéhyde ou une cétone.
\item Écrire les formules semi-développées et les noms des quatre aldéhydes et des trois cétones de formule $\ce{C5H10O}$.
\item Exploiter le spectre : quelle famille l'indice A désigne-t-il ? Justifier par la présence et l'absence de bandes.
\item Utiliser les indices B et C pour identifier la molécule. Écrire son nom.
\item Calculer sa masse molaire.
\item Pourquoi un seul indice ne suffisait-il pas ? Quel indice serait inutile si l'on savait que la chaîne principale compte
quatre atomes de carbone ?
\end{questions}
\end{blocEdition}

\editionPage{defi-ex-9}{D09 · Deux procédés pour l'ibuprofène}{Recherche et rédaction}{Avancé · 45 à 60 min · Chapitre 9\hfill\hyperref[defi-cor-9-1]{Corrigé p.~\pageref*{defi-cor-9-1}}}
\begin{blocEdition}{defi-ex-9}
L'ibuprofène, $\ce{C13H18O2}$, est un anti-inflammatoire produit à des milliers de tonnes par an. Le premier procédé industriel,
mis au point dans les années 1960, comptait six étapes ; un procédé des années 1990 n'en compte plus que trois.
\par\smallskip\noindent\textbf{L'économie d'atomes.} C'est le quotient de la masse molaire du produit par la somme des masses molaires
des réactifs, en tenant compte des nombres stœchiométriques : elle indique la part des atomes des réactifs qui se
retrouvent dans le produit. Elle vaut \qty{40}{\percent} pour le procédé en six étapes. Le procédé en trois étapes utilise, au total,
une molécule de chacun des réactifs suivants : isobutylbenzène $\ce{C10H14}$, anhydride éthanoïque $\ce{C4H6O3}$, dihydrogène $\ce{H2}$ et
monoxyde de carbone $\ce{CO}$.
\par\smallskip\noindent Masses molaires (\unit{g.mol^{-1}}) : H 1,0 ; C 12,0 ; O 16,0.
\begin{questions}
\item Dans une synthèse en plusieurs étapes, le produit de chaque étape est le réactif de la suivante. Montrer que le
rendement global est le produit des rendements des étapes.
\item On admet que chaque étape a un rendement de \qty{90}{\percent}. Calculer le rendement global de chaque procédé.
\item Calculer la masse molaire de l'ibuprofène, puis l'économie d'atomes du procédé en trois étapes.
\item Pour \qty{1}{kg} d'ibuprofène, quelle masse de « déchets » chaque procédé produit-il au minimum, si tous les rendements
valaient \qty{100}{\percent} ?
\item En tenant compte des rendements de la question 2, quelle masse de réactifs faut-il pour produire \qty{1}{kg} d'ibuprofène
avec chaque procédé ?
\item Écrire un programme Python qui affiche le rendement global d'une synthèse de 3, 6 et 10 étapes de rendement
\qty{90}{\percent}, et conclure sur l'intérêt des synthèses courtes pour une chimie plus durable.
\end{questions}
\end{blocEdition}

\editionPage{defi-ex-10}{D10 · Le meilleur carburant ?}{Recherche et rédaction}{Avancé · 45 à 60 min · Chapitre 10\hfill\hyperref[defi-cor-10-1]{Corrigé p.~\pageref*{defi-cor-10-1}}}
\begin{blocEdition}{defi-ex-10}
Pour un avion, une fusée ou un camion, un bon carburant doit libérer beaucoup d'énergie pour une faible masse, mais
aussi pour un faible volume de réservoir. On compare quatre carburants stockés à l'état liquide.
\begin{center}
\begin{tabular}{lcccc}
\toprule
Carburant & dihydrogène & méthane & kérosène ($\ce{C12H26}$) & éthanol \\
\midrule
$E_\mathrm{r}$ (\unit{kJ.mol^{-1}}) & $-242$ & $-802$ & -- & $-1235$ \\
Pouvoir calorifique (\unit{MJ.kg^{-1}}) & ? & ? & 43,0 & ? \\
Masse volumique du liquide (\unit{kg.L^{-1}}) & 0,071 & 0,42 & 0,80 & 0,789 \\
\bottomrule
\end{tabular}
\end{center}
Masses molaires (\unit{g.mol^{-1}}) : H 1,0 ; C 12,0 ; O 16,0.
\begin{questions}
\item Calculer le pouvoir calorifique massique du dihydrogène, du méthane et de l'éthanol.
\item Calculer, pour chaque carburant, l'énergie libérée par litre de liquide. Classer les carburants selon chaque critère.
\item Quelle masse et quel volume de dihydrogène, puis de kérosène, faut-il emporter pour disposer de \qty{1}{GJ} ?
\item Écrire l'équation de combustion du kérosène, puis calculer la masse de dioxyde de carbone rejetée par mégajoule.
Comparer au méthane et à l'éthanol.
\item Les premiers étages des fusées brûlent souvent du kérosène, les étages supérieurs du dihydrogène. Proposer une
explication.
\item Écrire un programme Python qui calcule l'énergie par litre de chaque carburant et les classe par ordre décroissant.
\end{questions}
\end{blocEdition}

\editionPage{defi-ex-11}{D11 · Ce qui tient le noyau}{Raisonnement et ordres de grandeur}{Avancé · 45 à 60 min · Chapitre 11\hfill\hyperref[defi-cor-11-1]{Corrigé p.~\pageref*{defi-cor-11-1}}}
\begin{blocEdition}{defi-ex-11}
Un noyau atomique contient des protons, tous chargés positivement, serrés les uns contre les autres. Ils devraient se
repousser violemment. Pourquoi le noyau ne vole-t-il pas en éclats ? On étudie deux protons du noyau d'hélium.
\schema{defi-11}
Données : $e=\qty{1,60e-19}{C}$ ; $m_\mathrm{p}=\qty{1,67e-27}{kg}$ ; $k=\qty{9,0e9}{N.m^2.C^{-2}}$ ; $G=\qty{6,67e-11}{N.m^2.kg^{-2}}$ ; $g=\qty{9,81}{N.kg^{-1}}$.
\begin{questions}
\item Calculer la valeur de la force électrostatique entre les deux protons. Est-elle attractive ou répulsive ?
\item Quelle masse faudrait-il suspendre à un fil, sur Terre, pour que son poids ait cette valeur ?
\item Calculer la force de gravitation entre les deux protons et la comparer au poids d'un proton à la surface de la Terre.
\item Calculer le rapport des deux forces. La gravitation peut-elle assurer la cohésion du noyau ?
\item À quelle distance la répulsion entre deux protons ne vaudrait-elle plus que \qty{1}{N} ?
\item Écrire un programme Python qui calcule ce rapport pour $d=\qty{2,0e-15}{m}$, \qty{1,0e-10}{m} et \qty{1}{m}. Interpréter.
\item Conclure : quelles propriétés doit posséder l'interaction qui maintient le noyau ?
\end{questions}
\end{blocEdition}

\editionPage{defi-ex-12}{D12 · Les hémisphères de Magdebourg}{Raisonnement et programmation}{Avancé · 45 à 60 min · Chapitre 12\hfill\hyperref[defi-cor-12-1]{Corrigé p.~\pageref*{defi-cor-12-1}}}
\begin{blocEdition}{defi-ex-12}
En 1654, Otto von Guericke, bourgmestre de Magdebourg, accole deux demi-sphères de cuivre et pompe l'air qu'elles
contiennent. Devant l'empereur, deux attelages de huit chevaux ne parviennent pas à les séparer.
\schema{defi-12}
Données : $R=\qty{25}{cm}$ ; volume intérieur $V=\qty{65,4}{L}$ ; cylindre de la pompe : $v=\qty{2,0}{L}$ ; $P_\mathrm{atm}=\qty{1,013}{bar}=\qty{1,013e5}{Pa}$. On admet que la
force pressante résultante sur une demi-sphère est celle qui s'exercerait sur un disque de rayon $R$ : $F=(P_\mathrm{atm}-P)\,\pi\,R^2$.
\begin{questions}
\item Calculer la force nécessaire pour séparer les demi-sphères si le vide intérieur était parfait.
\item Quelle force chaque cheval devrait-il exercer ? On pourrait remplacer l'un des attelages par un arbre : justifier.
\item À chaque coup de pompe, l'air intérieur se détend dans le cylindre (volume total $V+v$) à température constante, puis l'air du
cylindre est rejeté. Montrer que la pression est multipliée par $\dfrac{V}{V+v}$ à chaque coup.
\item Écrire un programme Python qui compte les coups nécessaires pour descendre à \qty{0,10}{bar}.
\item Calculer la force nécessaire pour séparer les demi-sphères avec cette pression résiduelle.
\item Une ventouse de rayon \qty{5,0}{cm} est plaquée sur une vitre. Quelle force maximale peut-elle supporter ? Pourquoi est-ce bien
moins en pratique ?
\end{questions}
\end{blocEdition}

\editionPage{defi-ex-13}{D13 · Freiner à temps}{Raisonnement et programmation}{Avancé · 45 à 60 min · Chapitre 13\hfill\hyperref[defi-cor-13-1]{Corrigé p.~\pageref*{defi-cor-13-1}}}
\begin{blocEdition}{defi-ex-13}
Un conducteur aperçoit un obstacle. Avant de freiner, il lui faut environ \qty{1,0}{s} pour réagir ; ensuite, la route exerce sur les
pneus une force de freinage qui peut atteindre \qty{80}{\percent} du poids de la voiture sur route sèche et \qty{50}{\percent} sur route mouillée.
\schema{defi-13}
Données : $g=\qty{9,81}{N.kg^{-1}}$ ; on étudie une voiture de \qty{1200}{kg} ; la force de freinage est supposée constante.
\begin{questions}
\item Quelle est la nature du mouvement pendant le temps de réaction ? Calculer la distance parcourue à 50, 90 et \qty{130}{km/h}.
\item Pendant le freinage sur route sèche, calculer la somme des forces, puis la variation de vitesse toutes les \qty{0,01}{s}. Dépend-elle
de la masse de la voiture ?
\item Écrire un programme Python qui simule le freinage pas à pas ($\Delta t=\qty{0,01}{s}$) et affiche la distance de freinage depuis \qty{130}{km/h}.
\item Le programme donne, sur route sèche puis mouillée : \qty{12,4}{m} et \qty{19,7}{m} à \qty{50}{km/h}, \qty{39,9}{m} et \qty{63,8}{m} à \qty{90}{km/h},
\qty{83,3}{m} et \qty{133,1}{m} à \qty{130}{km/h}. Calculer les six distances d'arrêt.
\item Comparer les distances de freinage à \qty{50}{km/h} et à \qty{100}{km/h} (\qty{49,3}{m} sur route sèche). Que se passe-t-il quand la vitesse double ?
\item Sur autoroute, le code de la route impose de rester au moins deux secondes derrière le véhicule qui précède. Cette distance
suffit-elle pour s'arrêter ? Pourquoi la règle est-elle pourtant raisonnable ?
\end{questions}
\end{blocEdition}

\editionPage{defi-ex-14}{D14 · Ne pas faire sauter les plombs}{Raisonnement et programmation}{Avancé · 45 à 60 min · Chapitre 14\hfill\hyperref[defi-cor-14-1]{Corrigé p.~\pageref*{defi-cor-14-1}}}
\begin{blocEdition}{defi-ex-14}
Un dimanche midi, le four, la plaque de cuisson et la bouilloire fonctionnent en même temps : tout s'éteint. On étudie
l'installation électrique de la cuisine, alimentée sous \qty{230}{V} ; on admet que $P=U\,I$ pour chaque appareil.
\schema{defi-14}
\begin{questions}
\item Calculer l'intensité appelée par chaque appareil. Quel circuit a le calibre le plus élevé ? Pourquoi ?
\item La bouilloire et le lave-linge sont branchés sur le même circuit de prises. Peuvent-ils fonctionner ensemble ?
\item Quelle intensité maximale l'abonnement de \qty{6}{kVA} autorise-t-il ? Expliquer la coupure du dimanche midi.
\item Écrire un programme Python qui passe en revue toutes les combinaisons d'appareils en marche et compte celles qui
dépassent \qty{6000}{W}.
\item Pour brancher la bouilloire loin d'une prise, on utilise une rallonge de \qty{25}{m} (soit \qty{50}{m} de fil de cuivre, aller et retour) de
section \qty{1,5}{mm^2}. Sa résistance vaut $R=\rho\,\dfrac{L}{S}$ avec $\rho=\qty{1,7e-8}{\ohm.m}$. Calculer la puissance dissipée dans la rallonge.
\item Pourquoi recommande-t-on de dérouler entièrement une rallonge enrouleur avant de l'utiliser à forte puissance ?
\end{questions}
\end{blocEdition}

\editionPage{defi-ex-15}{D15 · Rouler moins vite pour consommer moins}{Raisonnement et programmation}{Avancé · 45 à 60 min · Chapitre 15\hfill\hyperref[defi-cor-15-1]{Corrigé p.~\pageref*{defi-cor-15-1}}}
\begin{blocEdition}{defi-ex-15}
Abaisser la vitesse maximale sur autoroute de 130 à \qty{110}{km/h} permettrait d'économiser du carburant. On cherche à l'estimer pour
une voiture qui roule à vitesse constante sur une autoroute horizontale.
\par\smallskip
\textbf{Document 1. Les forces résistantes.} La résistance au roulement des pneus vaut environ \qty{150}{N}. Les frottements de l'air valent
$f=\dfrac{1}{2}\,\rho\,C_x S\,v^2$, avec $\rho=\qty{1,2}{kg.m^{-3}}$ et $C_x S=\qty{0,65}{m^2}$ pour cette voiture ($v$ en \unit{m.s^{-1}}).\par\smallskip
\textbf{Document 2. Le moteur.} Un litre d'essence libère \qty{35}{MJ} en brûlant ; le moteur en convertit \qty{30}{\percent} en travail mécanique.
Brûler un litre d'essence rejette environ \qty{2,3}{kg} de dioxyde de carbone.
\begin{questions}
\item À vitesse constante, que vaut la variation d'énergie cinétique de la voiture ? Que peut-on en déduire sur la somme des travaux des
forces qui s'exercent sur elle ?
\item Calculer la force des frottements de l'air à 90, 110 et \qty{130}{km/h}.
\item En déduire le travail que doit fournir le moteur sur \qty{100}{km} à \qty{130}{km/h}, puis le volume d'essence consommé.
\item Écrire un programme Python qui affiche la consommation pour 100 km à 90, 110 et \qty{130}{km/h}.
\item Calculer l'économie de carburant et de dioxyde de carbone pour 100 km en roulant à \qty{110}{km/h} au lieu de \qty{130}{km/h}.
\item Combien de temps perd-on sur 100 km ? Discuter.
\end{questions}
\end{blocEdition}

\editionPage{defi-ex-16}{D16 · Stocker l'électricité dans un lac}{Recherche et calcul}{Avancé · 45 à 60 min · Chapitre 16\hfill\hyperref[defi-cor-16-1]{Corrigé p.~\pageref*{defi-cor-16-1}}}
\begin{blocEdition}{defi-ex-16}
L'électricité se stocke mal. Les stations de transfert d'énergie par pompage (STEP) utilisent deux lacs séparés par un grand
dénivelé : on y stocke de l'énergie sous forme d'énergie potentielle de l'eau.
\schema{defi-16}
Données : $g=\qty{9,81}{N.kg^{-1}}$ ; masse volumique de l'eau \qty{1000}{kg.m^{-3}} ; rendement du pompage \num{0,85}, du turbinage \num{0,90} ; $\qty{1}{GWh}=\qty{3,6e12}{J}$.
\begin{questions}
\item Calculer l'énergie potentielle gagnée par l'eau du réservoir supérieur lorsqu'il est rempli (on prend le dénivelé de \qty{900}{m} pour toute l'eau).
\item Calculer l'énergie électrique restituée en turbinant toute cette eau, puis l'énergie consommée pour la pomper. En déduire le
rendement global du stockage.
\item Représenter le bilan énergétique d'un cycle pompage-turbinage.
\item Écrire un programme Python qui affiche l'énergie restituée, en MWh, pour des volumes de $10^6$, $10^7$ et \qty{e8}{m^3}.
\item La centrale peut fournir \qty{1800}{MW}. Quel débit d'eau faut-il ? Combien de temps peut-elle fonctionner à pleine puissance ?
\item Comparer l'énergie stockée à celle des batteries de voitures électriques (\qty{52}{kWh} chacune). Conclure.
\end{questions}
\end{blocEdition}

\editionPage{defi-ex-17}{D17 · Localiser l'épicentre d'un séisme}{Raisonnement et programmation}{Avancé · 45 à 60 min · Chapitre 17\hfill\hyperref[defi-cor-17-1]{Corrigé p.~\pageref*{defi-cor-17-1}}}
\begin{blocEdition}{defi-ex-17}
Trois stations sismiques A, B et C enregistrent le même séisme. Chacune mesure le retard entre l'arrivée des ondes S et celle
des ondes P. On cherche la position de l'épicentre sur la carte.
\schema{defi-17}
Données : célérités des ondes P \qty{6,0}{km.s^{-1}} et des ondes S \qty{3,5}{km.s^{-1}} ; on assimile la région à un plan.
\begin{questions}
\item Montrer que la distance $d$ entre une station et l'épicentre est liée au retard $\Delta t$ par $d=\dfrac{\Delta t}{k}$, avec $k=\dfrac{1}{3{,}5}-\dfrac{1}{6{,}0}$ (en s/km).
\item Calculer la distance de l'épicentre à chacune des trois stations.
\item Tracer sur la carte les trois cercles correspondants et en déduire la position de l'épicentre.
\item Écrire un programme Python qui teste tous les points de la carte, de kilomètre en kilomètre, et retient celui dont les distances aux
trois stations s'accordent le mieux avec les mesures.
\item À quelle date, après le début du séisme, les ondes P sont-elles arrivées en A ?
\item Pourquoi deux stations ne suffiraient-elles pas ?
\end{questions}
\end{blocEdition}

\editionPage{defi-ex-18}{D18 · Photographier la Lune}{Raisonnement et programmation}{Avancé · 45 à 60 min · Chapitre 18\hfill\hyperref[defi-cor-18-1]{Corrigé p.~\pageref*{defi-cor-18-1}}}
\begin{blocEdition}{defi-ex-18}
Sur la photo d'une pleine Lune prise au smartphone, la Lune n'est qu'un petit point blanc. Pourquoi ? Quel objectif faudrait-il
pour qu'elle remplisse la photo ?
\schema{defi-18}
Données : diamètre de la Lune \qty{3474}{km} ; distance Terre-Lune \qty{384400}{km} ; capteur d'un appareil reflex \qty{24}{mm} × \qty{36}{mm} ; objectif d'un
smartphone de focale \qty{4}{mm}, dont les pixels mesurent \qty{1,4}{\micro m}.
\begin{questions}
\item Montrer que l'image de la Lune se forme pratiquement dans le plan focal de l'objectif.
\item À l'aide du rayon passant par le centre optique, montrer que le diamètre de l'image vaut $d'=f'\times\alpha$, où $\alpha=\dfrac{D}{d}$ est le
diamètre apparent de la Lune (en radians). Calculer $\alpha$.
\item Calculer le diamètre de l'image avec un téléobjectif de \qty{300}{mm}. Quelle fraction de la hauteur du capteur occupe-t-elle ?
\item Calculer le diamètre de l'image avec le smartphone, en micromètres puis en pixels. Conclure.
\item Écrire un programme Python qui affiche le diamètre de l'image pour des focales de 4, 50, 300 et \qty{1200}{mm}.
\item Quelle focale faudrait-il pour que la Lune remplisse la hauteur du capteur ? Pourquoi les astronomes utilisent-ils des télescopes ?
\end{questions}
\end{blocEdition}

\editionPage{defi-ex-19}{D19 · Pourquoi le ciel est-il bleu ?}{Raisonnement et programmation}{Avancé · 45 à 60 min · Chapitre 19\hfill\hyperref[defi-cor-19-1]{Corrigé p.~\pageref*{defi-cor-19-1}}}
\begin{blocEdition}{defi-ex-19}
La lumière du Soleil est blanche ; pourtant le ciel est bleu le jour et le Soleil couchant paraît rouge. L'explication tient à la façon
dont les molécules de l'air diffusent la lumière.
\schema{defi-19}
Données : l'intensité de la lumière diffusée par les molécules de l'air est proportionnelle à $\dfrac{1}{\lambda^4}$ (diffusion de Rayleigh). Couleurs
repères : violet \qty{400}{nm}, bleu \qty{450}{nm}, vert \qty{550}{nm}, rouge \qty{700}{nm}.
\begin{questions}
\item Combien de fois le bleu est-il plus diffusé que le rouge ? Et que le vert ?
\item Écrire un programme Python qui affiche, pour les longueurs d'onde de 400 à \qty{700}{nm} (pas de \qty{50}{nm}), la diffusion relative à celle du rouge
(\qty{700}{nm}).
\item Expliquer pourquoi le ciel nous paraît bleu dans toutes les directions.
\item Le violet est encore plus diffusé que le bleu. Pourquoi le ciel ne paraît-il pas violet ? (Penser au spectre du Soleil et aux cônes.)
\item Au coucher, le Soleil est à \ang{5} au-dessus de l'horizon. Montrer que la lumière traverse environ 11 fois plus d'air qu'à midi. Expliquer
la couleur du Soleil couchant.
\item Les gouttelettes des nuages, bien plus grosses que les molécules, diffusent toutes les couleurs de la même façon. Expliquer la
couleur des nuages.
\end{questions}
\end{blocEdition}

\editionPage{defi-ex-20}{D20 · De quoi les étoiles sont-elles faites ?}{Raisonnement et programmation}{Avancé · 45 à 60 min · Chapitre 20\hfill\hyperref[defi-cor-20-1]{Corrigé p.~\pageref*{defi-cor-20-1}}}
\begin{blocEdition}{defi-ex-20}
Personne n'ira jamais prélever un échantillon d'étoile. Pourtant, on connaît la composition de leur atmosphère : il suffit d'analyser
leur lumière. Un astronome a mesuré sept raies sombres dans le spectre d'une étoile.
\schema{defi-20}
Données : longueurs d'onde (en nm) de quelques raies de référence, mesurées au laboratoire : hydrogène H 656,3 ; 486,1 ; 434,0 ;
410,2 — sodium Na 589,0 ; 589,6 — magnésium Mg 517,3 ; 518,4 — fer Fe 527,0 ; 438,4 — calcium Ca 393,4 ; 396,8. Niveaux de l'hydrogène :
$E_n=-\dfrac{\qty{13,6}{eV}}{n^2}$ ; $h=\qty{6,63e-34}{J.s}$ ; $c=\qty{3,00e8}{m.s^{-1}}$ ; $\qty{1}{eV}=\qty{1,60e-19}{J}$.
\begin{questions}
\item S'agit-il d'un spectre d'émission ou d'absorption ? Expliquer l'origine des raies sombres.
\item Montrer que la raie à \qty{434,0}{nm} correspond à une transition entre les niveaux $n=2$ et $n=5$ de l'hydrogène. D'où vient le petit
écart ?
\item Écrire un programme Python qui compare chaque raie mesurée aux raies de référence (écart inférieur à \qty{0,5}{nm}) et affiche la liste
des éléments présents, sans doublon.
\item Calculer l'énergie des photons absorbés par le calcium à \qty{393,4}{nm}. Dans quel domaine se trouve cette raie ?
\item À température ambiante, l'hydrogène n'absorbe aucune radiation visible. Que peut-on en déduire sur l'atmosphère de l'étoile ?
\item En 1868, on observe dans le spectre du Soleil une raie jaune à \qty{587,6}{nm} qui ne correspond à aucun élément connu. Le programme
l'aurait-il attribuée au sodium ? Que pouvait-on conclure ?
\end{questions}
\end{blocEdition}

\editionPage{defi-ex-21}{D21 · Combien de mesures faut-il ?}{Raisonnement et programmation}{Avancé · 45 à 60 min · Chapitre 21\hfill\hyperref[defi-cor-21-1]{Corrigé p.~\pageref*{defi-cor-21-1}}}
\begin{blocEdition}{defi-ex-21}
Pour connaître plus précisément un volume à l'équivalence, on peut répéter le titrage. Mais combien de fois ? Plutôt que de passer des
semaines au laboratoire, on simule les mesures : chaque titrage « virtuel » donne une valeur tirée au hasard autour de \qty{12,20}{mL},
avec une dispersion caractérisée par un écart-type de \qty{0,08}{mL}.
\begin{python}
import random
from statistics import mean, stdev
from math import sqrt
random.seed(2026)                    # memes tirages a chaque execution
for n in [5, 20, 100, 500]:
    mesures = [random.gauss(12.20, 0.08) for i in range(n)]
    s = stdev(mesures)
    print(n, round(mean(mesures), 3), round(s, 3), round(s / sqrt(n), 4))
\end{python}
Données : \code{random.gauss(m, e)} renvoie un nombre tiré au hasard, réparti « en cloche » autour de \code{m}, avec un écart-type \code{e}.
\begin{questions}
\item Que représente chaque nombre de la liste \code{mesures} ? Quel est l'intérêt de l'instruction \code{random.seed(2026)} ?
\item Exécuter le programme. Que représentent les quatre nombres affichés sur chaque ligne ? Comment évoluent la moyenne, l'écart-type $s$ et
l'incertitude-type $u$ quand $n$ augmente ?
\item Montrer que multiplier le nombre de mesures par 100 divise l'incertitude-type par 10.
\item Combien de titrages faudrait-il pour obtenir $u=\qty{0,001}{mL}$ ? Combien de temps faudrait-il, à raison d'un titrage de \qty{15}{min} ?
\item On remplace \code{12.20} par \code{12.30} pour simuler une burette dont le zéro est mal réglé. Que devient le résultat quand $n$ augmente ?
\item Proposer des moyens plus efficaces que la répétition pour améliorer la précision d'un titrage.
\end{questions}
\end{blocEdition}
